\section{Consultas representativas}

El conjunto contiene siete consultas funcionales y dos pruebas transversales.
Cada caso define sus parámetros, el resultado esperado y un orden estable con
códigos de negocio.

\begin{table}[ht]
\centering
\small
\begin{tabularx}{\textwidth}{>{\raggedright\arraybackslash}p{0.06\linewidth}>{\raggedright\arraybackslash}p{0.30\linewidth}Y}
\toprule
Caso & Pregunta o control & Patrón y resultado comprobado \\
\midrule
C1 & Condición comercial vigente. & Filtra un intervalo semiabierto por cliente, producto y fecha; devuelve una fila aplicable. \\
C2 & Pedidos y entregas que requieren atención. & Combina cuatro entidades con \texttt{JOIN}/\texttt{LEFT JOIN}; incluye demoras y parciales, y excluye el pedido cancelado. \\
C3 & Importe neto por categoría. & Usa \texttt{GROUP BY} sobre cantidad, precio y descuento; devuelve totales y conteos exactos. \\
C4 & Productos principales por segmento. & Aplica \texttt{DENSE\_RANK}; el código de producto sólo estabiliza la presentación de empates. \\
C5 & Clientes sin pedidos históricos. & Usa \texttt{NOT EXISTS}; devuelve el único cliente reservado para ese escenario. \\
C6 & Búsqueda vectorial top-k. & Ordena por distancia coseno y fragmento; el fragmento 10 ocupa la primera posición. \\
C7 & Recuperación híbrida autorizada. & Combina similitud, vigencia, modelo activo y RLS; Comercial/Compras obtiene primero el fragmento 18 y Operaciones no ve su clase. \\
C8 & Matriz de autorización. & Recorre las 15 combinaciones perfil--clase y confirma siete denegaciones por ausencia de permiso. \\
C9 & Trazabilidad e inmutabilidad. & Correlaciona actor, consulta, respuesta, evidencia y auditoría; las modificaciones operativas de eventos son rechazadas. \\
\bottomrule
\end{tabularx}
\caption{Siete consultas y dos pruebas transversales.}
\end{table}

C1--C7 producen los siete pares \texttt{consulta}--\texttt{respuesta}. C8 y C9
reutilizan esas interacciones para comprobar aislamiento y trazabilidad; no
crean otras conversaciones. Por eso hay nueve casos SQL y siete interacciones
simuladas.
